% =====================================================================
%  CONJURA CONJECTURE STATEMENT
%  Simultaneous inversion of a one-way function and a random permutation
%  in the presence of Simon's collision finder.
%  Requires conjura-conjecture.cls in the same directory.
% =====================================================================
\documentclass{conjura-conjecture}

\runninghead{SIMULTANEOUS INVERSION WITH A COLLISION FINDER}

\newcommand{\Fs}{\mathsf{F}}          % the one-way function oracle
\newcommand{\RP}{\mathsf{RP}}         % the random permutation
\newcommand{\Coll}{\mathsf{Coll}}     % Simon's collision finder
\newcommand{\Oc}{\mathsf{O}}
\newcommand{\bits}{\{0,1\}}
\newcommand{\Perm}{\mathrm{Perm}}
\newcommand{\negl}{\mathrm{negl}}

\begin{document}

\cjkicker{CONJURA \textperiodcentered{} OPEN PROBLEM}
\cjtitle{Simultaneous Inversion of a One-Way Function and a Random
Permutation With a Collision Finder}
\cjsubtitle{The strong amplification conjecture behind Simon's oracle}
\cjstatus{Statement: AI-written from Geoffroy Couteau, Pooya Farshim and
  Mohammad Mahmoody, \emph{Black-Box Uselessness: Composing Separations
  in Cryptography}, IACR Cryptology ePrint Archive 2021, not yet checked
  by a human. Open: the paper's weaker version without the collision
  finder is settled only in the case where the hard function is itself a
  random oracle and is open for an arbitrary one-way function, so the
  version stated here, with the collision finder, is open in every case;
  what the paper does establish is the two consequences, that this
  conjecture implies both the distributional and the class-reduction
  relaxations of its main helpfulness conjecture.}
\cjcategory{\emph{Impossibility Results \& Black-Box Separations}.}

\vspace{0.4em}

\begin{informalconjecture}
Fix any function family that is hard to invert, in the sense that no
efficient attacker succeeds on a random input with probability better
than some bound. Now independently sample a random permutation, and hand
the attacker a further oracle that, given any circuit built out of the
permutation and out of recursive calls to that same further oracle,
returns a random collision for it. The attacker is handed two challenges
at once, the image of a random input under the hard function and the
image of a random input under the permutation, and must recover both
inputs. The conjecture is that its success probability is at most the
hardness bound for the function, multiplied by the probability of
guessing a permutation preimage outright, up to a polynomial factor. In
words: the two inversion tasks are so independent that the attacker can
do no better than solve them separately, and the collision-finding
oracle does not help. The paper's weaker version of this statement, in
which the collision finder is dropped, is known when the hard function
is itself a random oracle but not when it is an arbitrary one-way
function; the version with the collision finder is stated here as a
conjecture and is not claimed to be known in any case. Beyond its own
interest as a statement about random permutations, the conjecture
matters because the collision-finding oracle is exactly Simon's, the one
used to separate collision-resistant hashing from one-way functions, and
because a known construction turns any pair of simultaneously
hard-to-invert functions into a collision-resistant hash function.
Proving it therefore settles two weakenings of the question of whether
one-way functions can help build collision resistance.
\end{informalconjecture}

\section{The setting}\label{sec:setting}

Simon~\cite{Sim98} separated collision-resistant hash functions from
one-way functions by exhibiting a random permutation $\RP$ together with
a collision finder $\Coll^{\RP}$, relative to which one-way permutations
exist but collision resistance does not. It is standard~\cite{IR89} that
a random permutation cannot be inverted with probability better than
$\poly(\lambda)/2^{\lambda}$ for a measure-one choice of the
permutation, and Simon's analysis shows that this survives the presence
of the collision finder, and indeed the presence of any independent
auxiliary oracle.

The conjecture asks for the corresponding \emph{simultaneous} statement
against an arbitrary independent hard function. The paper's motivation
is a route to its main helpfulness conjecture. Holmgren and
Lombardi~\cite{HL18} show how to build collision-resistant hash
functions in a black-box way from exponentially secure one-way product
functions, i.e.\ tuples of functions no efficient adversary can invert
simultaneously on random inputs with probability better than
$\negl(n)/2^{n}$; the paper notes in a footnote that their result as
stated requires the two components to be equal and injective, with an
extra-loss black-box reduction from an arbitrary pair. Taking one
component to be an arbitrary one-way function $\Fs$ and the other to be
$\RP$, the conjecture supplies exactly the product hardness the
construction needs, and so yields an auxiliary primitive, namely a
random permutation packaged with a collision finder, from which no
collision-resistant hash function can be built alone, by Simon's
separation, but from which one can be built once any one-way function is
added. The paper proves that this route works for two relaxed notions of
black-box reduction: one where reductions need only succeed for a
measure-one set of implementations of the auxiliary primitive, and one
where they are class reductions in the sense of Shaltiel~\cite{Sha20},
restricted to the efficient one-way functions. The paper also flags a
weaker version of the conjecture in which the collision finder is
dropped, notes that it is known when $\Fs$ is itself a random oracle but
not for an arbitrary one-way function, and says its study might be of
interest beyond black-box helpfulness, for instance in hardness
amplification. A footnote suggests an even weaker stepping stone: an
$\varepsilon$-secure one-way permutation together with a plain random
oracle. A second, independent route goes through the backdoored random
oracle model of Bauer, Farshim and Mazaheri~\cite{BFM18}, whose leakage
oracle can implement Simon's collision finder, conditional on a
communication-complexity conjecture related to set intersection, and
yielding in Section~6.4 the conclusion that a single backdoored random
oracle, though black-box separated from collision-resistant hashing, is
black-box helpful for it.

\section{Notation and parameters}\label{sec:notation}

Throughout, $\lambda\in\mathbb{N}$ is the security parameter and PPT
abbreviates probabilistic polynomial time. An \emph{oracle-aided PPT
machine} $A^{\Oc_1,\dots,\Oc_k}$ is a PPT machine with oracle access to
$\Oc_1,\dots,\Oc_k$ that runs in probabilistic polynomial time for every
choice of oracles, each oracle call counting as one step. We write
$x\sample S$ for a uniform draw from a finite set $S$.

Let $\Perm_\lambda$ denote the set of permutations of $\bits^{\lambda}$.
We write
$\Fs=\{\Fs_\lambda\colon\bits^{\lambda}\to\bits^{*}\}_{
\lambda\in\mathbb{N}}$ for a family of functions regarded as a single
oracle, $\RP=\{\RP_\lambda\}_{\lambda\in\mathbb{N}}$ with
$\RP_\lambda\in\Perm_\lambda$ for a permutation family regarded as a
single oracle, and $\Coll^{\RP}$ for the collision finder attached to
$\RP$. The function $\varepsilon\colon\mathbb{N}\to[0,1]$ measures the
hardness of inverting $\Fs$, and $q$ denotes a positive polynomial. We
write $\pi$ and $\Coll^{\pi}$ when referring to the same two objects as
they are named in Simon's original separation.

The parameters of the statement are:
\begin{itemize}
\item $\lambda$, the security parameter, also the input length of both
  challenge functions.
\item $\varepsilon$, the inversion hardness of the one-way function
  oracle; universally quantified over functions from the naturals to the
  unit interval.
\item $\Fs$, an arbitrary $\varepsilon$-secure one-way function family,
  independent of the permutation.
\item $\RP$, a uniformly random permutation family on bit strings of
  each length.
\item $\Coll^{\RP}$, Simon's collision finder for circuits built from
  the permutation and from recursive calls to itself.
\item $A$, the attacker, an oracle-aided PPT machine given all three
  oracles and both challenges.
\item $q$, the polynomial slack in the conclusion, allowed to depend on
  the attacker.
\end{itemize}

\section{The conjecture}\label{sec:conjectures}

\begin{definition}[$\varepsilon$-secure one-way function
oracle]\label{def:owf}
Let $\varepsilon\colon\mathbb{N}\to[0,1]$. A family
$\Fs=\{\Fs_\lambda\colon\bits^{\lambda}\to\bits^{*}\}_{
\lambda\in\mathbb{N}}$ is \emph{$\varepsilon$-secure one-way} if for
every oracle-aided PPT machine $B$ and all sufficiently large
$\lambda\in\mathbb{N}$,
\begin{align*}
  \Pr_{x\sample\bits^{\lambda}}\big[\Fs_{\lambda}\big(&
    B^{\Fs}(1^{\lambda},\Fs_{\lambda}(x))\big)\\
  &=\Fs_{\lambda}(x)\big]\ \le\ \varepsilon(\lambda).
\end{align*}
\end{definition}

\begin{definition}[Random permutation with Simon's collision
finder]\label{def:coll}
The \emph{random permutation}
$\RP=\{\RP_\lambda\}_{\lambda\in\mathbb{N}}$ is obtained by drawing
$\RP_\lambda\sample\Perm_\lambda$ independently for every $\lambda$; the
associated measure on families is the product of the uniform measures.

The \emph{collision finder} $\Coll^{\RP}$ is a randomised oracle defined
by induction on the nesting depth of $\Coll$-gates. On input a circuit
$C$ with $\RP$-gates and $\Coll^{\RP}$-gates computing a function
$\bits^{m}\to\bits^{m'}$, it returns a pair $(u,u')$ obtained by drawing
$u\sample\bits^{m}$ and then
$u'\sample\{v\in\bits^{m}: C(v)=C(u)\}$, so that $C(u)=C(u')$. Fixing,
for every circuit, the coins used in this sampling yields a
deterministic collision finder for $\RP$, and \emph{with probability $1$
over the choice of $\Coll^{\RP}$} refers to the induced measure on these
deterministic choices. This is the oracle written $\Coll^{\pi}$ for a
random permutation $\pi$ in Simon's separation.
\end{definition}

\begin{conjecture}[Simultaneous inversion of $\Fs$ and $\RP$ given
$\Coll^{\RP}$]\label{conj:main}
Let $\varepsilon\colon\mathbb{N}\to[0,1]$. For every family
$\Fs=\{\Fs_\lambda\}_{\lambda\in\mathbb{N}}$ that is
$\varepsilon$-secure one-way (Definition~\ref{def:owf}) and every
oracle-aided PPT machine $A$, there exists a polynomial $q$ such that,
with probability $1$ over the choice of the random permutation $\RP$ and
of the collision finder $\Coll^{\RP}$ (Definition~\ref{def:coll}), for
all sufficiently large $\lambda\in\mathbb{N}$,
\begin{multline*}
  \Pr\big[(x_1,x_2)=A^{\Fs,\RP,\Coll^{\RP}}\big(1^{\lambda},\\
    \Fs_{\lambda}(x_1),\ \RP_{\lambda}(x_2)\big)\big]\\
  \le\ \varepsilon(\lambda)\cdot\frac{q(\lambda)}{2^{\lambda}},
\end{multline*}
where the probability is taken over
$(x_1,x_2)\sample(\bits^{\lambda})^{2}$ and the coins of $A$.
\end{conjecture}

\begin{remark}[Independence of the two challenges]\label{rem:indep}
Note that $\Fs$ is fixed before $\RP$ and $\Coll^{\RP}$ are sampled, so
it is independent of them.
\end{remark}

\begin{remark}[Openness]\label{rem:open}
The paper does not prove the conjecture: ``We did not manage to prove
this conjecture, and leave it as an interesting open problem which might
be of independent interest'' (p.~8). Of the weak version without the
collision finder it records that ``The above conjecture is known in the
case where F is itself a random oracle, but not when it is an arbitrary
one-way function'' (p.~31), and it introduces the version stated here
with ``In our context, though, we need an even stronger version, which
we still deem plausible, where the random permutation comes with a
collision finder Coll'' (p.~31). What the paper does prove are two
consequences: Theorem~6.8, that the conjecture yields distributional
black-box helpfulness of one-way functions for collision-resistant hash
functions, and Theorem~6.12, that it yields class helpfulness over the
class of efficient one-way functions, the latter via a Borel--Cantelli
argument extracting a single auxiliary implementation good for all
countably many efficient one-way functions. Both proofs are given as
sketches, and both instantiate the pair of simultaneously hard functions
as the given one-way function together with the random permutation,
invoking Simon's analysis to see that the permutation stays
exponentially one-way in the presence of the collision finder and of any
independent oracle.
\end{remark}

\begin{remark}[Caveats for a reviewer]\label{rem:caveats}
The printed conjecture writes the first challenge as $\mathsf{P}(x_1)$
rather than $\Fs(x_1)$ in both the weak and the strong version. It is
written here as $\Fs(x_1)$, which is what the surrounding prose and the
proof of Theorem~6.8 require, but a reviewer should confirm this
reading. Second, the printed display carries no quantifier over
$\lambda$ and leaves $\poly(\lambda)$ unquantified; this has been read
as ``there is a polynomial $q$, depending on the attacker, such that the
bound holds for all sufficiently large $\lambda$'', and a reading with
``for all $\lambda$'' would be formally stronger. Third, the conjecture
as printed demands exact recovery of both preimages,
$(x_1,x_2)=A(\dots)$, whereas the statement derived from it in the proof
of Theorem~6.8 is the weaker any-preimage form; the paper's
exact-recovery form has been kept. Fourth, the definition of the
collision finder above spells out the recursion by nesting depth and the
induced measure on deterministic finders, details the paper leaves
to~\cite{Sim98}; the intended object is Simon's, and any discrepancy
should be resolved in Simon's favour. Finally, the post-2021 literature
has not been checked exhaustively for a proof or a counterexample.
\end{remark}

\section{Bibliography}

\begin{conjurabibliography}{99}

\bibitem[BFM18]{BFM18}
Balthazar Bauer, Pooya Farshim, and Sogol Mazaheri.
\emph{Combiners for backdoored random oracles}.
CRYPTO 2018, Part II, LNCS 10992, pages 272--302, Springer, 2018.

\bibitem[HL18]{HL18}
Justin Holmgren and Alex Lombardi.
\emph{Cryptographic hashing from strong one-way functions (or: One-way
product functions and their applications)}.
59th Annual Symposium on Foundations of Computer Science, pages
850--858, IEEE Computer Society Press, 2018.

\bibitem[IR89]{IR89}
Russell Impagliazzo and Steven Rudich.
\emph{Limits on the provable consequences of one-way permutations}.
21st Annual ACM Symposium on Theory of Computing, pages 44--61, ACM
Press, 1989.

\bibitem[Sha20]{Sha20}
Ronen Shaltiel.
\emph{Is it possible to improve Yao's XOR lemma using reductions that
exploit the efficiency of their oracle?}
APPROX/RANDOM 2020, Schloss Dagstuhl--Leibniz-Zentrum fur Informatik,
2020.

\bibitem[Sim98]{Sim98}
Daniel R. Simon.
\emph{Finding collisions on a one-way street: Can secure hash functions
be based on general assumptions?}
EUROCRYPT'98, LNCS 1403, pages 334--345, Springer, 1998.

\end{conjurabibliography}

\end{document}
